% ----------------------
\subsection{Section 107 (Between circa 1 March and circa 4 May 1835)}
% ----------------------
	\label{DC107}
	
	% -----
	\subsubsection{Relation to section 68}
	% -----
	
		\hyperref[DC68-14]{DC 68:14--21} includes passages about the priesthood, ``literal descendants'' of Aaron, and so on; none of that material is in the original version. It was added later, after DC 107 came along; see the \hyperref[DC68-HB]{historical background} to section 68.
	
	% -----
	\subsubsection{Historical Background}
	% -----
		\label{DC107-HB}
		
		% --
		\paragraph{JSP}
		% --
		
			``In early 1835, twelve men were called to serve as apostles in the Church of the Latter Day Saints, and numerous elders were appointed as members of the Seventy. In March 1835, JS and the Twelve Apostles decided that because of `the many pressing requests from the eastern churches,' the apostles would conduct a series of conferences in the eastern United States. These conferences were held `for the purpose of regulateing all things necessary' for the welfare of the branches of the church in those areas. It appears that sometime before the Twelve departed on 4 May 1835 to begin holding these conferences, JS presented them with this instruction, which outlined information about priesthood and church governance. The document indicated that there were two major divisions of priesthood in the church---the Melchizedek priesthood and the Aaronic priesthood---and explained the responsibilities of the Twelve Apostles, the Seventy, bishops, and other officers in the church. This instruction was apparently meant to aid the Twelve in their regulation of the eastern branches, while also providing detailed information about the responsibilities of priesthood officers. Like the earlier `Articles and Covenants' of the church, the instruction became an important document establishing significant doctrines on the governing bodies of the church and on the priesthood itself.
			
			``The Instruction on Priesthood compiled information from previous JS revelations. Much of the instruction came from a revelation JS had dictated in November 1831.%
				\footnote{%
					% \label{}%
					This is referring to a revelation given on \hyperref[EMS-1831-JS-11-NOV-1831-B]{11 November 1831}.
					}
			According to some observers, JS dictated other parts of the instruction as an 1835 revelation. Some of the instruction also elaborated on ideas first presented in a September 1832 revelation, indicating that JS was gaining new understanding into concepts of priesthood and leadership.%
				\footnote{%
					% \label{}%
					This is a reference to \hyperref[DC84]{DC 84}.
					}
			Oliver Cowdery, who assisted in calling the Twelve and giving them their `charge' as apostles, was probably involved in the instruction's preparation; Brigham Young later remembered JS spending `two hours laboring with Elder Cowdery to get him to write' what Young called a `Revelation on Priesthood'---probably this instruction.%
				\footnote{%
					% \label{}%
					The JSP note here says, ``School of the Prophets Provo Records, 15 Apr. 1868, 5--6.''
					}
			Indeed, when the instruction was published in the 1835 edition of the Doctrine and Covenants, it was designated as an instruction `On Priesthood.' Church members saw the instruction as coming from God. Heber C. Kimball, for example, referred to the instruction as something `the Lord' had `bestow[ed] upon' the Saints through `Brother Joseph' and explained that the Twelve Apostles `praised the Lord' for its contents.
			
			``Exactly when in 1835 the instruction appeared in its complete form is unclear. Although it may have been prepared directly after the Twelve's calling on 14 February 1835, it was probably presented in written form no earlier than 1 March. At a meeting of the Twelve on 27 February 1835, JS asked the apostles to discuss the question, `What importance is attached to the callings of these twelve apostles differrent from the other callings and offices of the chu[r]ch,' suggesting that he had not at that point instructed them on this matter. JS explained during this meeting that it was `all important that the twelve should understand the power and authority of the priesthoods, for without this knowledge they can do nothing to profit.' JS also proclaimed that the Twelve were `called to a travelling high council to preside over all the churches of the saints among the gentiles when there is no presidency established.' The Twelve were `to travel and preach among the Gentiles until the Lord shall command them to go to the Jews.' They were also to hold `the keys of this ministry--- to unlock the door of the kingdom of heaven unto all nations and preach the Gospel unto every creation.' Similar explanations appear in the featured instruction. In addition, the instruction contains detailed information about the office of Seventy, which was first instituted in the church at meetings held on 28 February and 1 March 1835. It is unlikely that the instruction was completed before this office was formally introduced into the church.
			
			``Heber C. Kimball remembered that JS dictated the instruction on priesthood at a meeting of the Twelve one evening after the apostles had been called and most of them had been ordained. Kimball's account places the meeting sometime before the first week of April 1835. A later JS history placed this dictation on 28 March 1835 and referenced a document recounting a meeting of the Twelve where, after a period of `general confession,' they made a written request that JS obtain for them a revelation of God's `mind and will concerning our duty the coming season.' The history states that JS then dictated this instruction. However, this appears to be based on a misdating by the history's compilers since the request to JS, which was copied into Minute Book 1, is dated 28 March 1836, not 1835. In fact, JS was not in Kirtland, Ohio, on 28 March 1835. On that day, five members of the Twelve---William E. McLellin, Lyman or Luke Johnson, David W. Patten, Parley P. Pratt, and John F. Boynton---were with JS in Huntsburgh, Ohio, preaching and baptizing. Another five apostles were presumably in Kirtland. Huntsburgh is about seventeen miles from Kirtland, and according to McLellin's diary, JS, and presumably the apostles with him, did not leave Huntsburgh until 30 March.
			
			``It is possible that JS formally presented the instruction to the Twelve on 26 April 1835. The minutes of a meeting held on that date state, `This day, pursuant to previous appointment, the Twelve Apostles and the Seventy (a part of whom had already been chosen,) assembled in the temple (altho' unfinished.) with a numerous concourse of people in order to receive our charge and instructions from President Joseph Smith Jun relative to our mission and duties.' The minutes, however, do not provide more detail on the instruction JS gave. Another possibility is that the featured instruction was presented in a 28 April 1835 meeting of the Twelve. At that meeting, a motion was carried that each apostle `forgive one another every wrong that has existed among us,' which may have been the `general confession' mentioned in the later JS history. In any case, this instruction was likely given to the Twelve before they left Kirtland on 4 May 1835, and it was certainly prepared by 26 May 1835, when William W. Phelps sent his wife, Sally Waterman Phelps, printed copies of `the Six first forms of the Doctrine and Covenants,' which included the instruction.
			
			``Although not covering all operations of church administration, the instruction captures information about the coordination of church governance between bishops, the presidency of the high priesthood, and newer offices such as the high council, the Twelve Apostles, and the Seventy. It also gives specific names for the two priesthoods in the church---Melchizedek for the higher priesthood and Aaronic for the lower---and explains why those two names were applied to these priesthoods.%
				\footnote{%
					% \label{}%
					JSP note: ``Earlier revelations had associated Melchizedek with the higher priesthood and Aaron with the lesser priesthood. The Book of Mormon noted that Melchizedek was a high priest, and JS's Bible revision explained that `Melchisedec was ordained a priest after the order of the Son of God.' JS and Sidney Rigdon's account of a February 1832 vision of the afterlife explained that those in the highest kingdom of glory were `priests of the most high after the order of Melchesadeck which was after the order of Enoch which was after the order of the only begotten son.' A September 1832 revelation stated that Abraham received the higher priesthood from Melchizedek, whereas the lesser priesthood was given to Aaron and his seed. While at an April 1834 meeting, JS referred to the `priesthood of Aaron' as something distinct from the higher priesthood. When JS ordained Oliver Cowdery an assistant president in the presidency of the high priesthood, JS referred to the priesthood as being `after the order of Melchizedek.' (Book of Mormon, 1830 ed., 260 [Alma 13:14]; New Testament Revision 2, p. 139 [second numbering] [Joseph Smith Translation, Hebrews 7:3]; Vision, 16 Feb. 1832 [D\&C 76:56--57]; Revelation, 22--23 Sept. 1832 [D\&C 84:14, 18]; Minutes and Discourse, 21 Apr. 1834; Account of Meetings, Revelation, and Blessing, 5--6 Dec. 1834; see also Old Testament Revision 1, pp. 33--34 [Joseph Smith Translation, Genesis 14:26--33].)''
					} 
			The instruction explains that the presidency of the high priesthood has the right to officiate in all offices of the church and further illuminates the role of the bishop in the church, explaining that the bishopric is the presidency of the Aaronic priesthood. Essentially, this document updated the Articles and Covenants and other revelations on church administration and codified the roles of the new offices.
			
			``No manuscript copies of the complete instruction are extant. The featured version was published in the 1835 edition of the Doctrine and Covenants, which was issued in September of that year. The volume includes two sections: one containing lectures on the doctrine of the church; and the other containing the church's `Covenants and Commandments,' or JS's revelations. This instruction is the third document in the `Covenants and Commandments' portion, directly following a revelatory preface and the foundational Articles and Covenants of the church. It is likely that at least some members of the Twelve took copies of the instruction with them to the eastern United States to use in their conferences. Minutes of a conference held in Westfield, New York, in May, for example, mention the Twelve providing `much explanation' about `the nature and principles of church government,' items that the instruction addresses. An August 1835 letter from Kirtland church leaders to the Twelve also contains a reference that suggests the Twelve had an advanced copy of the published version of the instruction at that time.''
						
	% -----
	\subsubsection{Versions}
	% -----
		\label{DC107-V}
		
		See the \href{https://drive.google.com/file/d/1goAvNz8jS4R8slYfMG_db55Ty2z_vmgK/view?usp=sharing}{sections comparison} document.
	
		\begin{compactdesc}
			\item[JSP] \hyperref[EMS-1835-JS-CIR-1-MAR-4-MAY-1835]{Between circa 1 March and circa 4 May 1835}
			\item[BC] ---
			\item[DC 1835] Section 3, 82--89
			\item[DC 1844] Section 3, 100--111
			\item[TS] ---
			\item[MS] 1:1 (May 1840), 13--17
			\item[PGP 1851] 49--52
		\end{compactdesc}

	% -----
	\subsubsection{Section 107:1} % *DC107-1 
	% -----
		\label{DC107-1}

			\footnote{%
				% \label{}%
				Note that in the 1835 and 1844 DC versions of this revelation, there is a huge heading at the beginning which reads, ``ON PRIESTHOOD.''
				}%
		THERE are, in the church,%
			\footnote{%
				% \label{}%
				President Openshaw was the first to call my attention to the importance of this qualifier, ``in the church.'' It appears to mean, ``in the church as it is organized on this earth.'' See \hyperref[DC107-8]{verses 8--9}, among other places. As far as I can tell, there are no manuscript sources of these verses in the first part of the revelation, so there's nothing to check against.
				}
		two priesthoods, namely, the Melchizedek and Aaronic, including the
			\footnote{%
				% \label{}%
				This is one of three uses of the term ``Levitical Priesthood'' in the scriptures. It occurs on verse 6 below, and at \hyperref[HEBREWS-7-11]{Hebrews 7:11}.
				}%
		Levitical Priesthood.

		\begin{framed}
		\scriptsize
			\begin{compactdesc}
				\item[JSP] 1 There are, in the church, two priesthoods, namely: the Melchizedek, and the Aaronic, including the Levitical priesthood.
				% \item[BC] 
				% \item[]
			\end{compactdesc}
		\end{framed}

	% -----
	\subsubsection{Section 107:2} % *DC107-2 
	% -----
		\label{DC107-2}

		Why the first is called the Melchizedek Priesthood is because Melchizedek was such a great high priest.%
			\footnote{%
				% \label{}%
				The JST of \hyperref[JST-OTR1-GENESIS-14]{Genesis 14} has a lot more information about Melchizedek; see section c but particularly sections d and e.
				}

		\begin{framed}
		\scriptsize
			\begin{compactdesc}
				\item[JSP] Why the first is called the Melchizedek priesthood, is because Melchizedek was such a great high priest:
				% \item[BC] 
				% \item[]
			\end{compactdesc}
		\end{framed}

	% -----
	\subsubsection{Section 107:3} % *DC107-3 
	% -----
		\label{DC107-3}

		Before his day it was called \textit{the Holy Priesthood, after the} \hyperref[ORDER]{\textit{Order}} \textit{of the Son of God}.%
			\footnote{%
				% \label{}%
				Compare \hyperref[DC84-18]{DC 84:18} and \hyperref[DC93-20]{DC 93:20}. For more on this ``order,'' see also \hyperref[DC107-40]{DC 107:40--57}; JST OTR1 Genesis 14 \hyperlink{JST-OTR1-GENESIS-14-d}{section d}; \hyperref[ALMA-13-1]{Alma 13:1--16}. For other related ideas and commentary, see \hyperref[JOHN-5-26]{John 5:26} and notes there, 
				}

		\begin{framed}
		\scriptsize
			\begin{compactdesc}
				\item[JSP] before his day it was called \textit{the holy priesthood, after the order of the Son of God;}
				% \item[BC] 
				% \item[]
			\end{compactdesc}
		\end{framed}

	% -----
	\subsubsection{Section 107:4} % *DC107-4 
	% -----
		\label{DC107-4}

		But out of respect or reverence to the name of the Supreme Being, to avoid the too frequent repetition of his name, they, the church, in ancient days, called that priesthood after Melchizedek, or the Melchizedek Priesthood.

		\begin{framed}
		\scriptsize
			\begin{compactdesc}
				\item[JSP] but out of respect or reverence to the name of the Supreme Being, to avoid the too frequent repetition of his name, they, the church, in ancient days, called that priesthood after Melchizedek, or the Melchizedek priesthood.
				% \item[BC] 
				% \item[]
			\end{compactdesc}
		\end{framed}

	% -----
	\subsubsection{Section 107:5} % *DC107-5 
	% -----
		\label{DC107-5}

		All other authorities or offices in the church are \hyperref[APPENDAGE]{appendages} to this priesthood.

		% \begin{framed}
		% \scriptsize
			% \begin{compactdesc}
				% \item[JSP] 
				% % \item[BC] 
				% % \item[]
			% \end{compactdesc}
		% \end{framed}

	% -----
	\subsubsection{Section 107:6} % *DC107-6 
	% -----
		\label{DC107-6}

		But there are two divisions or grand heads%
			\footnote{%
				% \label{}%
				Definition \#20 of the 1828 dictionary for ``head'' is ``Topic of discourse; chief point or subject; a summary; as the heads of a discourse or treatise.'' As an adjective, the word means ``chief, principal,'' so it could be a substantive adjective here. One way to understand this remark here is to think of the ``heads'' of the Melchizedek and Aaronic priesthoods as subsets of the whole system of laws that compose the priesthood generally; for human beings, these two subsets happen to be very important. See the JS 5 January 1841 remark at \hyperref[PRIESTHOOD-JS]{Priesthood}.
				}%
		---one is the Melchizedek Priesthood, and the other is the Aaronic or Levitical Priesthood.%
			\footnote{%
				% \label{}%
				Here the Aaronic priesthood seems to be equated with the Levitical priesthood; it may be that they can officiate in the same ordinances, but I think the difference has to do with the line of descent.
				}

		\begin{framed}
		\scriptsize
			\begin{compactdesc}
				\item[JSP] but there are two divisions, or grand heads---one is the Melchizedek priesthood, and the other is the Aaronic, or Levitical priesthood.
				% \item[BC] 
				% \item[]
			\end{compactdesc}
		\end{framed}

	% -----
	\subsubsection{Section 107:7} % *DC107-7 
	% -----
		\label{DC107-7}

		The office of an elder comes under the priesthood of Melchizedek.

		% \begin{framed}
		% \scriptsize
			% \begin{compactdesc}
				% \item[JSP] 
				% % \item[BC] 
				% % \item[]
			% \end{compactdesc}
		% \end{framed}

	% -----
	\subsubsection{Section 107:8} % *DC107-8 
	% -----
		\label{DC107-8}

		The Melchizedek Priesthood holds the right of presidency, and has power and authority over all the offices in the church in all ages of the world, to administer in spiritual things.

		% \begin{framed}
		% \scriptsize
			% \begin{compactdesc}
				% \item[JSP] 
				% % \item[BC] 
				% % \item[]
			% \end{compactdesc}
		% \end{framed}

	% -----
	\subsubsection{Section 107:9} % *DC107-9 
	% -----
		\label{DC107-9}

		The Presidency of the High Priesthood, after the order of Melchizedek, have a right to officiate in all the offices in the church.

		% \begin{framed}
		% \scriptsize
			% \begin{compactdesc}
				% \item[JSP] 
				% % \item[BC] 
				% % \item[]
			% \end{compactdesc}
		% \end{framed}

	% -----
	\subsubsection{Section 107:10} % *DC107-10 
	% -----
		\label{DC107-10}

		High priests after the order of the Melchizedek Priesthood have a right to officiate in their own standing, under the direction of the presidency, in administering spiritual things, and also in the office of an elder, priest (of the Levitical order), teacher, deacon, and member.

		% \begin{framed}
		% \scriptsize
			% \begin{compactdesc}
				% \item[JSP] 
				% % \item[BC] 
				% % \item[]
			% \end{compactdesc}
		% \end{framed}

	% -----
	\subsubsection{Section 107:11} % *DC107-11 
	% -----
		\label{DC107-11}

		An elder has a right to officiate in his stead when the high priest is not present.

		% \begin{framed}
		% \scriptsize
			% \begin{compactdesc}
				% \item[JSP] 
				% % \item[BC] 
				% % \item[]
			% \end{compactdesc}
		% \end{framed}

	% -----
	\subsubsection{Section 107:12} % *DC107-12 
	% -----
		\label{DC107-12}

		The high priest and elder are to administer in spiritual things, agreeable to the covenants and commandments of the church; and they have a right to officiate in all these offices of the church when there are no higher authorities present.

		% \begin{framed}
		% \scriptsize
			% \begin{compactdesc}
				% \item[JSP] 
				% % \item[BC] 
				% % \item[]
			% \end{compactdesc}
		% \end{framed}

	% -----
	\subsubsection{Section 107:13} % *DC107-13 
	% -----
		\label{DC107-13}

		The second priesthood is called the Priesthood of Aaron, because it was conferred upon Aaron and his seed, throughout all their generations.

		% \begin{framed}
		% \scriptsize
			% \begin{compactdesc}
				% \item[JSP] 
				% % \item[BC] 
				% % \item[]
			% \end{compactdesc}
		% \end{framed}

	% -----
	\subsubsection{Section 107:14} % *DC107-14 
	% -----
		\label{DC107-14}

		Why it is called the lesser priesthood is because it is an \hyperref[APPENDAGE]{appendage} to the greater, or the Melchizedek Priesthood, and has power in administering outward%
			\footnote{%
				% \label{DC107-14-outward}%
				See discussion of the phrase ``outward ordinances'' and definitions of ``outward'' \hyperref[ORDINANCE-PHRASES]{here}. The contrast with ``outward'' appears to be ``spiritual,'' as in ``spiritual blessings'' in \hyperref[DC107-18]{verse 18}.
				}
		 ordinances.

		% \begin{framed}
		% \scriptsize
			% \begin{compactdesc}
				% \item[JSP] 
				% % \item[BC] 
				% % \item[]
			% \end{compactdesc}
		% \end{framed}

	% -----
	\subsubsection{Section 107:15} % *DC107-15 
	% -----
		\label{DC107-15}

		The bishopric is the presidency of this priesthood, and holds the keys or authority of the same.

		% \begin{framed}
		% \scriptsize
			% \begin{compactdesc}
				% \item[JSP] 
				% % \item[BC] 
				% % \item[]
			% \end{compactdesc}
		% \end{framed}

	% -----
	\subsubsection{Section 107:16} % *DC107-16 
	% -----
		\label{DC107-16}

		No man has a legal right to this office, to hold the keys of this priesthood, except he be a literal descendant of Aaron.%
			\footnote{%
				% \label{}%
				The way this discussion goes, it also seems like it's referring to the Levitical priesthood, and maybe it would just restrict the line of descent to those of Levite descent.
				}

		% \begin{framed}
		% \scriptsize
			% \begin{compactdesc}
				% \item[JSP] 
				% % \item[BC] 
				% % \item[]
			% \end{compactdesc}
		% \end{framed}

	% -----
	\subsubsection{Section 107:17} % *DC107-17 
	% -----
		\label{DC107-17}

		But as a high priest of the Melchizedek Priesthood has authority to officiate in all the lesser offices, he may officiate in the office of bishop when no literal descendant of Aaron can be found, provided he is called and set apart and ordained unto this power by the hands of the Presidency of the Melchizedek Priesthood.

		% \begin{framed}
		% \scriptsize
			% \begin{compactdesc}
				% \item[JSP] 
				% % \item[BC] 
				% % \item[]
			% \end{compactdesc}
		% \end{framed}

	% -----
	\subsubsection{Section 107:18} % *DC107-18 
	% -----
		\label{DC107-18}

			\footnote{%
				% \label{}%
				I read this verse and the next as enumerating a list of things that constitute the ``power and authority'' of the Melchizedek priesthood. In particular, I take the one at the end of this verse, ``to hold the keys of all the spiritual blessings of the church,'' as the first item on this list, and then ``to have the privilege of receiving the mysteries of the kingdom of heaven'' as the second item, and so on, rather than seeeing the ``privilege'' and the ``heavens opened'' one as explications of what the spiritual blessings are.
				}%
		The power and authority of the higher, or Melchizedek Priesthood, is to hold the keys of all the spiritual blessings of the church---

		\begin{framed}
		\scriptsize
			\begin{compactdesc}
				\item[JSP] The power and authority of the higher or Melchizedek priesthood, is to hold the keys of all the spiritual blessings of the church---
				% \item[BC] 
				% \item[]
			\end{compactdesc}
		\end{framed}

	% -----
	\subsubsection{Section 107:19} % *DC107-19 
	% -----
		\label{DC107-19}

		To have the privilege of receiving the \hyperref[MYSTERY]{mysteries} of the kingdom of heaven, to have the heavens opened unto them,
			\footnote{%
				% \label{}%
				Compare \hyperref[HEBREWS-12-22]{Hebrews 12:22--24}.
				}%
		to commune with the general assembly and church of the \hyperref[FIRSTBORN]{Firstborn},%
			\footnote{%
				% \label{}%
				For ``church of the Firstborn,'' see \hyperref[CHURCH-PHRASES]{here}. Note that in the JSP (1835 DC) version, it's ``church of the first born.''
				}
		and to enjoy the communion and presence of God the Father, and Jesus the mediator of the new covenant.

		\begin{framed}
		\scriptsize
			\begin{compactdesc}
				\item[JSP] to have the privilege of receiving the mysteries of the kingdom of heaven---to have the heavens opened unto them---to commune with the general assembly and church of the first born, and to enjoy the communion and presence of God the Father, and Jesus the Mediator of the new covenant.
				% \item[BC] 
				% \item[]
			\end{compactdesc}
		\end{framed}

	% -----
	\subsubsection{Section 107:20} % *DC107-20 
	% -----
		\label{DC107-20}

		The power and authority of the lesser, or Aaronic Priesthood, is to hold the keys of the ministering of angels, and to administer in \hyperref[OUTWARD]{outward} ordinances, the letter%
			\footnote{%
				% \label{}%
				See \hyperref[2CORINTHIANS-3-6]{2 Corinthians 3:6}?
				} 
		of the gospel, the baptism of repentance for the remission of sins, agreeable to the covenants and commandments.%
			\footnote{%
				% \label{}%
				Compare \hyperref[DC13-1]{DC 13} and \hyperref[DC84-26]{DC 84:26--27}.
				}

		\begin{framed}
		\scriptsize
			\begin{compactdesc}
				\item[JSP] The power and authority of the lesser, or Aaronic priesthood, is, to hold the keys of the ministring of angels, and to administer in outward ordinances---the letter of the gospel---the baptism of repentance for the remission of sins, agreeably to the covenants and commandments.
				% \item[BC] 
				% \item[]
			\end{compactdesc}
		\end{framed}

	% -----
	\subsubsection{Section 107:21} % *DC107-21 
	% -----
		\label{DC107-21}

		Of necessity there are presidents, or presiding officers growing out of, or appointed of or from among those who are ordained to the several offices in these two priesthood.

		% \begin{framed}
		% \scriptsize
			% \begin{compactdesc}
				% \item[JSP] 
				% % \item[BC] 
				% % \item[]
			% \end{compactdesc}
		% \end{framed}

	% -----
	\subsubsection{Section 107:22} % *DC107-22 
	% -----
		\label{DC107-22}

		Of the Melchizedek Priesthood, three Presiding High Priests, chosen by the body, appointed and ordained to that office, and upheld by the confidence, faith, and prayer of the church, form a quorum of the Presidency of the Church.

		% \begin{framed}
		% \scriptsize
			% \begin{compactdesc}
				% \item[JSP] 
				% % \item[BC] 
				% % \item[]
			% \end{compactdesc}
		% \end{framed}

	% -----
	\subsubsection{Section 107:23} % *DC107-23 
	% -----
		\label{DC107-23}

		The twelve traveling councilors are called to be the Twelve Apostles, or special witnesses of the name of Christ in all the world---thus differing from other officers in the church in the duties of their calling.

		% \begin{framed}
		% \scriptsize
			% \begin{compactdesc}
				% \item[JSP] 
				% % \item[BC] 
				% % \item[]
			% \end{compactdesc}
		% \end{framed}

	% -----
	\subsubsection{Section 107:24} % *DC107-24 
	% -----
		\label{DC107-24}

		And they form a quorum, equal in authority and power to the three presidents previously mentioned.

		% \begin{framed}
		% \scriptsize
			% \begin{compactdesc}
				% \item[JSP] 
				% % \item[BC] 
				% % \item[]
			% \end{compactdesc}
		% \end{framed}

	% -----
	\subsubsection{Section 107:25} % *DC107-25 
	% -----
		\label{DC107-25}

		The Seventy are also called to preach the gospel, and to be especial witnesses unto the Gentiles and in all the world---thus differing from other officers in the church in the duties of their calling.

		% \begin{framed}
		% \scriptsize
			% \begin{compactdesc}
				% \item[JSP] 
				% % \item[BC] 
				% % \item[]
			% \end{compactdesc}
		% \end{framed}

	% -----
	\subsubsection{Section 107:26} % *DC107-26 
	% -----
		\label{DC107-26}

		And they form a quorum, equal in authority to that of the Twelve special witnesses or Apostles just named.

		% \begin{framed}
		% \scriptsize
			% \begin{compactdesc}
				% \item[JSP] 
				% % \item[BC] 
				% % \item[]
			% \end{compactdesc}
		% \end{framed}

	% -----
	\subsubsection{Section 107:27} % *DC107-27 
	% -----
		\label{DC107-27}

		And every decision made by either of these quorums must be by the unanimous voice of the same; that is, every member in each quorum must be agreed to its decisions, in order to make their decisions of the same power or validity one with the other---

		% \begin{framed}
		% \scriptsize
			% \begin{compactdesc}
				% \item[JSP] 
				% % \item[BC] 
				% % \item[]
			% \end{compactdesc}
		% \end{framed}

	% -----
	\subsubsection{Section 107:28} % *DC107-28 
	% -----
		\label{DC107-28}

		A majority may form a quorum when circumstances render it impossible to be otherwise---

		% \begin{framed}
		% \scriptsize
			% \begin{compactdesc}
				% \item[JSP] 
				% % \item[BC] 
				% % \item[]
			% \end{compactdesc}
		% \end{framed}

	% -----
	\subsubsection{Section 107:29} % *DC107-29 
	% -----
		\label{DC107-29}

		Unless this is the case, their decisions are not entitled to the same blessings which the decisions of a quorum of three presidents were anciently, who were ordained after the order of Melchizedek, and were righteous and holy men.

		% \begin{framed}
		% \scriptsize
			% \begin{compactdesc}
				% \item[JSP] 
				% % \item[BC] 
				% % \item[]
			% \end{compactdesc}
		% \end{framed}

	% -----
	\subsubsection{Section 107:30} % *DC107-30 
	% -----
		\label{DC107-30}

		The decisions of these quorums, or either of them, are to be made in all righteousness, in holiness, and lowliness of heart, meekness and long suffering, and in faith, and virtue, and knowledge, temperance, patience, godliness, brotherly kindness and charity;

		% \begin{framed}
		% \scriptsize
			% \begin{compactdesc}
				% \item[JSP] 
				% % \item[BC] 
				% % \item[]
			% \end{compactdesc}
		% \end{framed}

	% -----
	\subsubsection{Section 107:31} % *DC107-31 
	% -----
		\label{DC107-31}

		Because the promise is,
			\footnote{%
				% \label{}%
				The material in the previous verse and here is a reference to \hyperref[2PETER-1-5]{2 Peter 1:5--8}.
				}%
		if these things abound in them they shall not be unfruitful in the knowledge of the Lord.

		% \begin{framed}
		% \scriptsize
			% \begin{compactdesc}
				% \item[JSP] 
				% % \item[BC] 
				% % \item[]
			% \end{compactdesc}
		% \end{framed}

	% -----
	\subsubsection{Section 107:32} % *DC107-32 
	% -----
		\label{DC107-32}

		And in case that any decision of these quorums is made in unrighteousness, it may be brought before a general assembly of the several quorums, which constitute the spiritual authorities of the church; otherwise there can be no appeal from their decision.

		% \begin{framed}
		% \scriptsize
			% \begin{compactdesc}
				% \item[JSP] 
				% % \item[BC] 
				% % \item[]
			% \end{compactdesc}
		% \end{framed}

	% -----
	\subsubsection{Section 107:33} % *DC107-33 
	% -----
		\label{DC107-33}

		The Twelve are a Traveling Presiding High Council, to officiate in the name of the Lord, under the direction of the Presidency of the Church, agreeable to the institution of heaven; to build up the church, and regulate all the affairs of the same in all nations, first unto the Gentiles and secondly unto the Jews.

		% \begin{framed}
		% \scriptsize
			% \begin{compactdesc}
				% \item[JSP] 
				% % \item[BC] 
				% % \item[]
			% \end{compactdesc}
		% \end{framed}

	% -----
	\subsubsection{Section 107:34} % *DC107-34 
	% -----
		\label{DC107-34}

		The Seventy are to act in the name of the Lord, under the direction of the Twelve or the traveling high council, in building up the church and regulating all the affairs of the same in all nations, first unto the Gentiles and then to the Jews;

		% \begin{framed}
		% \scriptsize
			% \begin{compactdesc}
				% \item[JSP] 
				% % \item[BC] 
				% % \item[]
			% \end{compactdesc}
		% \end{framed}

	% -----
	\subsubsection{Section 107:35} % *DC107-35 
	% -----
		\label{DC107-35}

		The Twelve being sent out, holding the keys, to open the door by the proclamation of the gospel of Jesus Christ, and first unto the Gentiles and then unto the Jews.

		% \begin{framed}
		% \scriptsize
			% \begin{compactdesc}
				% \item[JSP] 
				% % \item[BC] 
				% % \item[]
			% \end{compactdesc}
		% \end{framed}

	% -----
	\subsubsection{Section 107:36} % *DC107-36 
	% -----
		\label{DC107-36}

		The standing high councils, at the stakes of Zion, form a quorum equal in authority in the affairs of the church, in all their decisions, to the quorum of the presidency, or to the traveling high council.

		% \begin{framed}
		% \scriptsize
			% \begin{compactdesc}
				% \item[JSP] 
				% % \item[BC] 
				% % \item[]
			% \end{compactdesc}
		% \end{framed}

	% -----
	\subsubsection{Section 107:37} % *DC107-37 
	% -----
		\label{DC107-37}

		The high council in Zion form a quorum equal in authority in the affairs of the church, in all their decisions, to the councils of the Twelve at the stakes of Zion.

		% \begin{framed}
		% \scriptsize
			% \begin{compactdesc}
				% \item[JSP] 
				% % \item[BC] 
				% % \item[]
			% \end{compactdesc}
		% \end{framed}

	% -----
	\subsubsection{Section 107:38} % *DC107-38 
	% -----
		\label{DC107-38}

		It is the duty of the traveling high council to call upon the Seventy, when they need assistance, to fill the several calls for preaching and administering the gospel, instead of any others.

		% \begin{framed}
		% \scriptsize
			% \begin{compactdesc}
				% \item[JSP] 
				% % \item[BC] 
				% % \item[]
			% \end{compactdesc}
		% \end{framed}

	% -----
	\subsubsection{Section 107:39} % *DC107-39 
	% -----
		\label{DC107-39}

		It is the duty of the Twelve, in all large branches of the church, to ordain evangelical ministers, as they shall be designated unto them by revelation---

		\begin{framed}
		\scriptsize
			\begin{compactdesc}
				\item[JSP] 17 It is the duty of the twelve in all large branches of the church, to ordain evangelical ministers, as they shall be designated unto them by revelation.
				% \item[BC] 
				% \item[]
			\end{compactdesc}
		\end{framed}

	% -----
	\subsubsection{Section 107:40} % *DC107-40 
	% -----
		\label{DC107-40}

			\footnote{%
				% \label{}%
				Notice how, in the JSP (1835 DC), the previous verse is a numbered section by itself and ends with a period. The way it is written now is very confusing, with the m-dash, because it makes it seem like verse 39 is somehow related to this one or is an introduction to it or something. That isn't right. Verse 40 here begins a new major division of this revelation (which goes until the end of verse 57).
				}%
		The order of this priesthood%
			\footnote{%
				% \label{}%
				I think we must be talking about the Melchizedek priesthood here, given that the text has just been naming offices and duties belonging to that priesthood.
				}
		was confirmed to be handed down from father to son, and rightly belongs to the literal descendants of the chosen seed,%
			\footnote{%
				% \label{}%
				Referring to Abraham here, maybe? This is why I think it's significant that Melchizedek was not a descendant of Abraham. But from verse 42 it could be Seth, and from the rest of the verses following that it looks like it should be Adam.
				}
		to whom the promises were made.

		% \begin{framed}
		% \scriptsize
			% \begin{compactdesc}
				% \item[JSP] 
				% % \item[BC] 
				% % \item[]
			% \end{compactdesc}
		% \end{framed}

	% -----
	\subsubsection{Section 107:41} % *DC107-41 
	% -----
		\label{DC107-41}

		This order was instituted in the days of Adam, and came down by lineage in the following manner:

		% \begin{framed}
		% \scriptsize
			% \begin{compactdesc}
				% \item[JSP] 
				% % \item[BC] 
				% % \item[]
			% \end{compactdesc}
		% \end{framed}

	% -----
	\subsubsection{Section 107:42} % *DC107-42 
	% -----
		\label{DC107-42}

		From Adam to Seth, who was ordained by Adam at the age of sixty-nine years, and was blessed by him three years previous to his (Adam's) death, and received the promise of God by his father, that his posterity should be the chosen of the Lord, and that they should be preserved unto the end of the earth;

		% \begin{framed}
		% \scriptsize
			% \begin{compactdesc}
				% \item[JSP] 
				% % \item[BC] 
				% % \item[]
			% \end{compactdesc}
		% \end{framed}

	% -----
	\subsubsection{Section 107:43} % *DC107-43 
	% -----
		\label{DC107-43}

		Because he (Seth) was a perfect man, and his likeness was the express likeness of his father, insomuch that he seemed to be like unto his father in all things, and could be distinguished from him only by his age.%
			\footnote{%
				% \label{}%
				I wonder if this verse could also be taken as a description of the relationship between Christ and the Father.
				}

		% \begin{framed}
		% \scriptsize
			% \begin{compactdesc}
				% \item[JSP] 
				% % \item[BC] 
				% % \item[]
			% \end{compactdesc}
		% \end{framed}

	% -----
	\subsubsection{Section 107:44} % *DC107-44 
	% -----
		\label{DC107-44}

		Enos was ordained at the age of one hundred and thirty-four years and four months, by the hand of Adam.

		% \begin{framed}
		% \scriptsize
			% \begin{compactdesc}
				% \item[JSP] 
				% % \item[BC] 
				% % \item[]
			% \end{compactdesc}
		% \end{framed}

	% -----
	\subsubsection{Section 107:45} % *DC107-45 
	% -----
		\label{DC107-45}

		God called upon Cainan in the wilderness in the fortieth year of his age; and he met Adam in journeying to the place Shedolamak. He was eighty-seven years old when he received his ordination.

		% \begin{framed}
		% \scriptsize
			% \begin{compactdesc}
				% \item[JSP] 
				% % \item[BC] 
				% % \item[]
			% \end{compactdesc}
		% \end{framed}

	% -----
	\subsubsection{Section 107:46} % *DC107-46 
	% -----
		\label{DC107-46}

		Mahalaleel was four hundred and ninety-six years and seven days old when he was ordained by the hand of Adam, who also blessed him.

		% \begin{framed}
		% \scriptsize
			% \begin{compactdesc}
				% \item[JSP] 
				% % \item[BC] 
				% % \item[]
			% \end{compactdesc}
		% \end{framed}

	% -----
	\subsubsection{Section 107:47} % *DC107-47 
	% -----
		\label{DC107-47}

		Jared was two hundred years old when he was ordained under the hand of Adam, who also blessed him.

		% \begin{framed}
		% \scriptsize
			% \begin{compactdesc}
				% \item[JSP] 
				% % \item[BC] 
				% % \item[]
			% \end{compactdesc}
		% \end{framed}

	% -----
	\subsubsection{Section 107:48} % *DC107-48 
	% -----
		\label{DC107-48}

		Enoch was twenty-five years old when he was ordained under the hand of Adam; and he was sixty-five and Adam blessed him.

		% \begin{framed}
		% \scriptsize
			% \begin{compactdesc}
				% \item[JSP] 
				% % \item[BC] 
				% % \item[]
			% \end{compactdesc}
		% \end{framed}

	% -----
	\subsubsection{Section 107:49} % *DC107-49 
	% -----
		\label{DC107-49}

		And he saw the Lord, and he walked with him, and was before his face continually; and he walked with God three hundred and sixty-five years, making him four hundred and thirty years old when he was translated.

		% \begin{framed}
		% \scriptsize
			% \begin{compactdesc}
				% \item[JSP] 
				% % \item[BC] 
				% % \item[]
			% \end{compactdesc}
		% \end{framed}

	% -----
	\subsubsection{Section 107:50} % *DC107-50 
	% -----
		\label{DC107-50}

		Methuselah was one hundred years old when he was ordained under the hand of Adam.

		% \begin{framed}
		% \scriptsize
			% \begin{compactdesc}
				% \item[JSP] 
				% % \item[BC] 
				% % \item[]
			% \end{compactdesc}
		% \end{framed}

	% -----
	\subsubsection{Section 107:51} % *DC107-51 
	% -----
		\label{DC107-51}

		Lamech was thirty-two years old when he was ordained under the hand of Seth.

		% \begin{framed}
		% \scriptsize
			% \begin{compactdesc}
				% \item[JSP] 
				% % \item[BC] 
				% % \item[]
			% \end{compactdesc}
		% \end{framed}

	% -----
	\subsubsection{Section 107:52} % *DC107-52 
	% -----
		\label{DC107-52}

		Noah was ten years old when he was ordained under the hand of Methuselah.

		% \begin{framed}
		% \scriptsize
			% \begin{compactdesc}
				% \item[JSP] 
				% % \item[BC] 
				% % \item[]
			% \end{compactdesc}
		% \end{framed}

	% -----
	\subsubsection{Section 107:53} % *DC107-53 
	% -----
		\label{DC107-53}

		Three years previous to the death of Adam, he called Seth, Enos, Cainan, Mahalaleel, Jared, Enoch, and Methuselah, who were all high priests, with the residue of his posterity who were righteous, into the valley of Adam-ondi-Ahman, and there bestowed upon them his last blessing.

		% \begin{framed}
		% \scriptsize
			% \begin{compactdesc}
				% \item[JSP] 
				% % \item[BC] 
				% % \item[]
			% \end{compactdesc}
		% \end{framed}

	% -----
	\subsubsection{Section 107:54} % *DC107-54 
	% -----
		\label{DC107-54}

		And the Lord appeared unto them, and they rose up and blessed Adam, and called him Michael, the prince, the archangel.

		% \begin{framed}
		% \scriptsize
			% \begin{compactdesc}
				% \item[JSP] 
				% % \item[BC] 
				% % \item[]
			% \end{compactdesc}
		% \end{framed}

	% -----
	\subsubsection{Section 107:55} % *DC107-55 
	% -----
		\label{DC107-55}

		And the Lord administered comfort unto Adam, and said unto him: I have set thee to be at the head; a multitude of nations shall come of thee,%
			\footnote{%
				% \label{}%
				Notice how these promises are very similar to the ones given to Abraham. There is someting important about the promise of posterity, especially when you think of JS and BY's later remarks about the constituents of a kingdom.
				} 
		and thou art a prince over them forever.

		% \begin{framed}
		% \scriptsize
			% \begin{compactdesc}
				% \item[JSP] 
				% % \item[BC] 
				% % \item[]
			% \end{compactdesc}
		% \end{framed}

	% -----
	\subsubsection{Section 107:56} % *DC107-56 
	% -----
		\label{DC107-56}

		And Adam stood up in the midst of the congregation; and, notwithstanding he was bowed down with age, being full of the Holy Ghost, predicted%
			\footnote{%
				% \label{}%
				This is the only use of ``predict'' anywhere in the scriptures.
				}
		whatsoever should befall his posterity unto the latest generation.

		% \begin{framed}
		% \scriptsize
			% \begin{compactdesc}
				% \item[JSP] 
				% % \item[BC] 
				% % \item[]
			% \end{compactdesc}
		% \end{framed}

	% -----
	\subsubsection{Section 107:57} % *DC107-57 
	% -----
		\label{DC107-57}

		These things were all written in the book of Enoch,%
			\footnote{%
				% \label{}%
				This appears to be the only reference in the scriptures to the ``book of Enoch.''
				} 
		and are to be testified of in due time.

		% \begin{framed}
		% \scriptsize
			% \begin{compactdesc}
				% \item[JSP] 
				% % \item[BC] 
				% % \item[]
			% \end{compactdesc}
		% \end{framed}

	% -----
	\subsubsection{Section 107:58} % *DC107-58 
	% -----
		\label{DC107-58}

		It is the duty of the Twelve, also, to ordain and set in order all the other officers of the church, agreeable to the revelation which says:

		% \begin{framed}
		% \scriptsize
			% \begin{compactdesc}
				% \item[JSP] 
				% % \item[BC] 
				% % \item[]
			% \end{compactdesc}
		% \end{framed}

	% -----
	\subsubsection{Section 107:59} % *DC107-59 
	% -----
		\label{DC107-59}

		To the church of Christ in the land of Zion, in addition to the church laws respecting church business---

		% \begin{framed}
		% \scriptsize
			% \begin{compactdesc}
				% \item[JSP] 
				% % \item[BC] 
				% % \item[]
			% \end{compactdesc}
		% \end{framed}

	% -----
	\subsubsection{Section 107:60} % *DC107-60 
	% -----
		\label{DC107-60}

		Verily, I say unto you, saith the Lord of Hosts, there must needs be presiding elders to preside over those who are of the office of an elder;

		% \begin{framed}
		% \scriptsize
			% \begin{compactdesc}
				% \item[JSP] 
				% % \item[BC] 
				% % \item[]
			% \end{compactdesc}
		% \end{framed}

	% -----
	\subsubsection{Section 107:61} % *DC107-61 
	% -----
		\label{DC107-61}

		And also priests to preside over those who are of the office of a priest;

		% \begin{framed}
		% \scriptsize
			% \begin{compactdesc}
				% \item[JSP] 
				% % \item[BC] 
				% % \item[]
			% \end{compactdesc}
		% \end{framed}

	% -----
	\subsubsection{Section 107:62} % *DC107-62 
	% -----
		\label{DC107-62}

		And also teachers to preside over those who are of the office of a teacher, in like manner, and also the deacons---

		% \begin{framed}
		% \scriptsize
			% \begin{compactdesc}
				% \item[JSP] 
				% % \item[BC] 
				% % \item[]
			% \end{compactdesc}
		% \end{framed}

	% -----
	\subsubsection{Section 107:63} % *DC107-63 
	% -----
		\label{DC107-63}

		Wherefore, from deacon to teacher, and from teacher to priest, and from priest to elder, severally as they are appointed, according to the covenants and commandments of the church.

		% \begin{framed}
		% \scriptsize
			% \begin{compactdesc}
				% \item[JSP] 
				% % \item[BC] 
				% % \item[]
			% \end{compactdesc}
		% \end{framed}

	% -----
	\subsubsection{Section 107:64} % *DC107-64 
	% -----
		\label{DC107-64}

		Then comes the High Priesthood, which is the greatest of all.

		% \begin{framed}
		% \scriptsize
			% \begin{compactdesc}
				% \item[JSP] 
				% % \item[BC] 
				% % \item[]
			% \end{compactdesc}
		% \end{framed}

	% -----
	\subsubsection{Section 107:65} % *DC107-65 
	% -----
		\label{DC107-65}

		Wherefore, it must needs be that one be appointed of the High Priesthood to preside over the priesthood, and he shall be called President of the High Priesthood of the Church;

		% \begin{framed}
		% \scriptsize
			% \begin{compactdesc}
				% \item[JSP] 
				% % \item[BC] 
				% % \item[]
			% \end{compactdesc}
		% \end{framed}

	% -----
	\subsubsection{Section 107:66} % *DC107-66 
	% -----
		\label{DC107-66}

		Or, in other words, the Presiding High Priest over the High Priesthood of the Church.

		% \begin{framed}
		% \scriptsize
			% \begin{compactdesc}
				% \item[JSP] 
				% % \item[BC] 
				% % \item[]
			% \end{compactdesc}
		% \end{framed}

	% -----
	\subsubsection{Section 107:67} % *DC107-67 
	% -----
		\label{DC107-67}

		From the same comes the administering of ordinances and blessings upon the church, by the laying on of the hands.

		% \begin{framed}
		% \scriptsize
			% \begin{compactdesc}
				% \item[JSP] 
				% % \item[BC] 
				% % \item[]
			% \end{compactdesc}
		% \end{framed}

	% -----
	\subsubsection{Section 107:68} % *DC107-68 
	% -----
		\label{DC107-68}

		Wherefore, the office of a bishop is not equal unto it; for the office of a bishop is in administering all temporal things;

		% \begin{framed}
		% \scriptsize
			% \begin{compactdesc}
				% \item[JSP] 
				% % \item[BC] 
				% % \item[]
			% \end{compactdesc}
		% \end{framed}

	% -----
	\subsubsection{Section 107:69} % *DC107-69 
	% -----
		\label{DC107-69}

		Nevertheless a bishop must be chosen from the High Priesthood, unless he is a literal descendant of Aaron;

		% \begin{framed}
		% \scriptsize
			% \begin{compactdesc}
				% \item[JSP] 
				% % \item[BC] 
				% % \item[]
			% \end{compactdesc}
		% \end{framed}

	% -----
	\subsubsection{Section 107:70} % *DC107-70 
	% -----
		\label{DC107-70}

		For unless he is a literal descendant of Aaron he cannot hold the keys of that priesthood.

		% \begin{framed}
		% \scriptsize
			% \begin{compactdesc}
				% \item[JSP] 
				% % \item[BC] 
				% % \item[]
			% \end{compactdesc}
		% \end{framed}

	% -----
	\subsubsection{Section 107:71} % *DC107-71 
	% -----
		\label{DC107-71}

		Nevertheless, a high priest, that is, after the order of Melchizedek, may be set apart unto the ministering of temporal things, having a knowledge of them by the Spirit of truth;

		% \begin{framed}
		% \scriptsize
			% \begin{compactdesc}
				% \item[JSP] 
				% % \item[BC] 
				% % \item[]
			% \end{compactdesc}
		% \end{framed}

	% -----
	\subsubsection{Section 107:72} % *DC107-72 
	% -----
		\label{DC107-72}

		And also to be a judge in Israel, to do the business of the church, to sit in judgment upon transgressors upon testimony as it shall be laid before him according to the laws, by the assistance of his counselors, whom he has chosen or will choose among the elders of the church.

		% \begin{framed}
		% \scriptsize
			% \begin{compactdesc}
				% \item[JSP] 
				% % \item[BC] 
				% % \item[]
			% \end{compactdesc}
		% \end{framed}

	% -----
	\subsubsection{Section 107:73} % *DC107-73 
	% -----
		\label{DC107-73}

		This is the duty of a bishop who is not a literal descendant of Aaron, but has been ordained to the High Priesthood after the order of Melchizedek.

		% \begin{framed}
		% \scriptsize
			% \begin{compactdesc}
				% \item[JSP] 
				% % \item[BC] 
				% % \item[]
			% \end{compactdesc}
		% \end{framed}

	% -----
	\subsubsection{Section 107:74} % *DC107-74 
	% -----
		\label{DC107-74}

		Thus shall he be a judge, even a common judge among the inhabitants of Zion, or in a stake of Zion, or in any branch of the church where he shall be set apart unto this ministry, until the borders of Zion are enlarged and it becomes necessary to have other bishops or judges in Zion or elsewhere.

		% \begin{framed}
		% \scriptsize
			% \begin{compactdesc}
				% \item[JSP] 
				% % \item[BC] 
				% % \item[]
			% \end{compactdesc}
		% \end{framed}

	% -----
	\subsubsection{Section 107:75} % *DC107-75 
	% -----
		\label{DC107-75}

		And inasmuch as there are other bishops appointed they shall act in the same office.

		% \begin{framed}
		% \scriptsize
			% \begin{compactdesc}
				% \item[JSP] 
				% % \item[BC] 
				% % \item[]
			% \end{compactdesc}
		% \end{framed}

	% -----
	\subsubsection{Section 107:76} % *DC107-76 
	% -----
		\label{DC107-76}

		But a literal descendant of Aaron has a legal right to the presidency of this priesthood, to the keys of this ministry, to act in the office of bishop independently, without counselors, except in a case where a President of the High Priesthood, after the order of Melchizedek, is tried, to sit as a judge in Israel.

		% \begin{framed}
		% \scriptsize
			% \begin{compactdesc}
				% \item[JSP] 
				% % \item[BC] 
				% % \item[]
			% \end{compactdesc}
		% \end{framed}

	% -----
	\subsubsection{Section 107:77} % *DC107-77 
	% -----
		\label{DC107-77}

		And the decision of either of these councils, agreeable to the commandment which says:

		% \begin{framed}
		% \scriptsize
			% \begin{compactdesc}
				% \item[JSP] 
				% % \item[BC] 
				% % \item[]
			% \end{compactdesc}
		% \end{framed}

	% -----
	\subsubsection{Section 107:78} % *DC107-78 
	% -----
		\label{DC107-78}

		Again, verily, I say unto you, the most important business of the church, and the most difficult cases of the church, inasmuch as there is not satisfaction upon the decision of the bishop or judges, it shall be handed over and carried up unto the council of the church, before the Presidency of the High Priesthood.

		% \begin{framed}
		% \scriptsize
			% \begin{compactdesc}
				% \item[JSP] 
				% % \item[BC] 
				% % \item[]
			% \end{compactdesc}
		% \end{framed}

	% -----
	\subsubsection{Section 107:79} % *DC107-79 
	% -----
		\label{DC107-79}

		And the Presidency of the council of the High Priesthood shall have power to call other high priests, even twelve, to assist as counselors; and thus the Presidency of the High Priesthood and its counselors shall have power to decide upon testimony according to the laws of the church.

		% \begin{framed}
		% \scriptsize
			% \begin{compactdesc}
				% \item[JSP] 
				% % \item[BC] 
				% % \item[]
			% \end{compactdesc}
		% \end{framed}

	% -----
	\subsubsection{Section 107:80} % *DC107-80 
	% -----
		\label{DC107-80}

		And after this decision it shall be had in remembrance no more before the Lord; for this is the highest council of the church of God, and a final decision upon controversies in spiritual matters.

		% \begin{framed}
		% \scriptsize
			% \begin{compactdesc}
				% \item[JSP] 
				% % \item[BC] 
				% % \item[]
			% \end{compactdesc}
		% \end{framed}

	% -----
	\subsubsection{Section 107:81} % *DC107-81 
	% -----
		\label{DC107-81}

		There is not any person belonging to the church who is exempt from this council of the church.

		% \begin{framed}
		% \scriptsize
			% \begin{compactdesc}
				% \item[JSP] 
				% % \item[BC] 
				% % \item[]
			% \end{compactdesc}
		% \end{framed}

	% -----
	\subsubsection{Section 107:82} % *DC107-82 
	% -----
		\label{DC107-82}

		And inasmuch as a President of the High Priesthood shall transgress, he shall be had in remembrance before the common council of the church, who shall be assisted by twelve counselors of the High Priesthood;

		% \begin{framed}
		% \scriptsize
			% \begin{compactdesc}
				% \item[JSP] 
				% % \item[BC] 
				% % \item[]
			% \end{compactdesc}
		% \end{framed}

	% -----
	\subsubsection{Section 107:83} % *DC107-83 
	% -----
		\label{DC107-83}

		And their decision upon his head shall be an end of controversy concerning him.

		% \begin{framed}
		% \scriptsize
			% \begin{compactdesc}
				% \item[JSP] 
				% % \item[BC] 
				% % \item[]
			% \end{compactdesc}
		% \end{framed}

	% -----
	\subsubsection{Section 107:84} % *DC107-84 
	% -----
		\label{DC107-84}

		Thus, none shall be exempted from the justice and the laws of God, that all things may be done in order and in solemnity before him, according to truth and righteousness.

		% \begin{framed}
		% \scriptsize
			% \begin{compactdesc}
				% \item[JSP] 
				% % \item[BC] 
				% % \item[]
			% \end{compactdesc}
		% \end{framed}

	% -----
	\subsubsection{Section 107:85} % *DC107-85 
	% -----
		\label{DC107-85}

		And again, verily I say unto you, the duty of a president over the office of a deacon is to preside over twelve deacons, to sit in council with them, and to teach them their duty, edifying one another, as it is given according to the covenants.

		% \begin{framed}
		% \scriptsize
			% \begin{compactdesc}
				% \item[JSP] 
				% % \item[BC] 
				% % \item[]
			% \end{compactdesc}
		% \end{framed}

	% -----
	\subsubsection{Section 107:86} % *DC107-86 
	% -----
		\label{DC107-86}

		And also the duty of the president over the office of the teachers is to preside over twenty-four of the teachers, and to sit in council with them, teaching them the duties of their office, as given in the covenants.

		% \begin{framed}
		% \scriptsize
			% \begin{compactdesc}
				% \item[JSP] 
				% % \item[BC] 
				% % \item[]
			% \end{compactdesc}
		% \end{framed}

	% -----
	\subsubsection{Section 107:87} % *DC107-87 
	% -----
		\label{DC107-87}

		Also the duty of the president over the Priesthood of Aaron is to preside over forty-eight priests, and sit in council with them, to teach them the duties of their office, as is given in the covenants---

		% \begin{framed}
		% \scriptsize
			% \begin{compactdesc}
				% \item[JSP] 
				% % \item[BC] 
				% % \item[]
			% \end{compactdesc}
		% \end{framed}

	% -----
	\subsubsection{Section 107:88} % *DC107-88 
	% -----
		\label{DC107-88}

		This president is to be a bishop; for this is one of the duties of this priesthood.

		% \begin{framed}
		% \scriptsize
			% \begin{compactdesc}
				% \item[JSP] 
				% % \item[BC] 
				% % \item[]
			% \end{compactdesc}
		% \end{framed}

	% -----
	\subsubsection{Section 107:89} % *DC107-89 
	% -----
		\label{DC107-89}

		Again, the duty of the president over the office of elders is to preside over ninety-six elders, and to sit in council with them, and to teach them according to the covenants.

		% \begin{framed}
		% \scriptsize
			% \begin{compactdesc}
				% \item[JSP] 
				% % \item[BC] 
				% % \item[]
			% \end{compactdesc}
		% \end{framed}

	% -----
	\subsubsection{Section 107:90} % *DC107-90 
	% -----
		\label{DC107-90}

		This presidency is a distinct one from that of the seventy, and is designed for those who do not travel into all the world.

		% \begin{framed}
		% \scriptsize
			% \begin{compactdesc}
				% \item[JSP] 
				% % \item[BC] 
				% % \item[]
			% \end{compactdesc}
		% \end{framed}

	% -----
	\subsubsection{Section 107:91} % *DC107-91 
	% -----
		\label{DC107-91}

		And again, the duty of the President of the office of the High Priesthood is to preside over the whole church, and to be like unto Moses---

		% \begin{framed}
		% \scriptsize
			% \begin{compactdesc}
				% \item[JSP] 
				% % \item[BC] 
				% % \item[]
			% \end{compactdesc}
		% \end{framed}

	% -----
	\subsubsection{Section 107:92} % *DC107-92 
	% -----
		\label{DC107-92}

		Behold, here is wisdom; yea, to be a seer, a revelator, a translator, and a prophet, having all the gifts of God which he bestows upon the head of the church.

		% \begin{framed}
		% \scriptsize
			% \begin{compactdesc}
				% \item[JSP] 
				% % \item[BC] 
				% % \item[]
			% \end{compactdesc}
		% \end{framed}

	% -----
	\subsubsection{Section 107:93} % *DC107-93 
	% -----
		\label{DC107-93}

		And it is according to the vision showing the order of the Seventy, that they should have seven presidents to preside over them, chosen out of the number of the seventy;

		% \begin{framed}
		% \scriptsize
			% \begin{compactdesc}
				% \item[JSP] 
				% % \item[BC] 
				% % \item[]
			% \end{compactdesc}
		% \end{framed}

	% -----
	\subsubsection{Section 107:94} % *DC107-94 
	% -----
		\label{DC107-94}

		And the seventh president of these presidents is to preside over the six;

		% \begin{framed}
		% \scriptsize
			% \begin{compactdesc}
				% \item[JSP] 
				% % \item[BC] 
				% % \item[]
			% \end{compactdesc}
		% \end{framed}

	% -----
	\subsubsection{Section 107:95} % *DC107-95 
	% -----
		\label{DC107-95}

		And these seven presidents are to choose other seventy besides the first seventy to whom they belong, and are to preside over them;

		% \begin{framed}
		% \scriptsize
			% \begin{compactdesc}
				% \item[JSP] 
				% % \item[BC] 
				% % \item[]
			% \end{compactdesc}
		% \end{framed}

	% -----
	\subsubsection{Section 107:96} % *DC107-96 
	% -----
		\label{DC107-96}

		And also other seventy, until seven times seventy, if the labor in the vineyard of necessity requires it.

		% \begin{framed}
		% \scriptsize
			% \begin{compactdesc}
				% \item[JSP] 
				% % \item[BC] 
				% % \item[]
			% \end{compactdesc}
		% \end{framed}

	% -----
	\subsubsection{Section 107:97} % *DC107-97 
	% -----
		\label{DC107-97}

		And these seventy are to be traveling ministers, unto the Gentiles first and also unto the Jews.

		% \begin{framed}
		% \scriptsize
			% \begin{compactdesc}
				% \item[JSP] 
				% % \item[BC] 
				% % \item[]
			% \end{compactdesc}
		% \end{framed}

	% -----
	\subsubsection{Section 107:98} % *DC107-98 
	% -----
		\label{DC107-98}

		Whereas other officers of the church, who belong not unto the Twelve, neither to the Seventy, are not under the responsibility to travel among all nations, but are to travel as their circumstances shall allow, notwithstanding they may hold as high and responsible offices in the church.

		% \begin{framed}
		% \scriptsize
			% \begin{compactdesc}
				% \item[JSP] 
				% % \item[BC] 
				% % \item[]
			% \end{compactdesc}
		% \end{framed}

	% -----
	\subsubsection{Section 107:99} % *DC107-99 
	% -----
		\label{DC107-99}

		Wherefore, now let every man learn his duty, and to act in the office in which he is appointed, in all diligence.

		% \begin{framed}
		% \scriptsize
			% \begin{compactdesc}
				% \item[JSP] 
				% % \item[BC] 
				% % \item[]
			% \end{compactdesc}
		% \end{framed}

	% -----
	\subsubsection{Section 107:100} % *DC107-100 
	% -----
		\label{DC107-100}

		He that is slothful shall not be counted worthy to stand, and he that learns not his duty and shows himself not approved shall not be counted worthy to stand. Even so. Amen.

		% \begin{framed}
		% \scriptsize
			% \begin{compactdesc}
				% \item[JSP] 
				% % \item[BC] 
				% % \item[]
			% \end{compactdesc}
		% \end{framed}